%% abstract.tex: abstract in PT and EN  (FEUP regulations)
%% -------------------------------------------------------
\chapter*{Resumo}

Este trabalho apresenta o projeto, implementação física, e verificação pós-layout de um transmissor \ac{CML} de 10~Gb/s na tecnologia CMOS GlobalFoundries 55~nm LPx, visando atender aos requisitos elétricos da norma USB~3.2 Gen~2. O transmissor utiliza uma arquitetura segmentada com \ac{FFE} de três derivações, com pré-cursor, cursor principal e pós-cursor independentemente programáveis. Os requisitos de \textit{preshoot} e \textit{de-emphasis} do USB são traduzidos em coeficientes para cada uma das três derivações através de valores de corrente, implementados por fontes de corrente digitalmente controláveis, enquanto a terminação é implementada com um banco de resistências digitalmente controlável, permitindo ajuste para compensar variações de processo, tensão, temperatura e variações do processo de fabricação dos resistências.

Uma metodologia projeto modular e abutável é usada para o \textit{layout}, desenvolvido utilizando posicionamento \textit{common-centroid}, roteamento simétrico, dispositivos \textit{dummy}, e \textit{guard rings} contínuos para reduzir desajustes e efeitos dependentes do \textit{layout}. O transmissor completo ocupa 176~$\mu$m $\times$ 82~$\mu$m. Simulações pré- e pós-extração são realizadas utilizando um canal de referência fornecido pela USB-IF com 21.6~dB de perda de inserção na frequência de Nyquist de 5~GHz. A extração parasítica afeta significativamente os tempos de transição de saída, o peso efetivo das derivações \ac{FFE}, e a configuração da terminação, sendo necessário \textit{tuning} das fontes de corrente e terminação para manter o desempenho em diferentes \textit{corners} operacionais.

Os resultados mostram que a arquitetura segmentada fornece uma faixa de ajuste suficiente para recuperar pontos de operação compatíveis com USB~3.2 Gen~2 após a implementação física, mas também demonstram que os códigos de equalização e terminação necessários podem se tornar dependentes do processo e das condições de operação. Em particular, a equalização pós-extração não pode ser representada por um mapeamento fixo único em todos os \textit{corners} PVT avaliados, o que destaca a importância dos mecanismos de calibração e a versatilidade da arquitetura proposta. O trabalho demonstra, portanto, uma metodologia completa de projeto para um transmissor configurável de 10 Gbps, desde o esquemático até ao \textit{layout}, identificando a calibração automática, caracterização completa de \textit{jitter} e \ac{BER}, e medição em silício como os passos principais necessários para validação experimental completa de conformidade.

\vspace{10mm}
\textbf{Palavras-chave:} Transmissor de alta velocidade, USB~3.2 Gen~2, Feed-Forward Equalization, Terminação, Projeto de circuitos integrados, Layout, Current Mode Logic.

\acresetall %% to reset the acronym usage

\chapter*{Abstract}
This dissertation presents the design, physical implementation, and post-layout verification of a 10~Gb/s \ac{CML} \ac{TX} targeting the USB~3.2 Gen~2 electrical requirements in GlobalFoundries 55~nm LPx CMOS technology. The transmitter employs a segmented design, with three-tap \ac{FFE}, composed of independently programmable pre-cursor, main-cursor, and post-cursor paths. The USB preshoot and de-emphasis requirements are translated into tap coefficients and current targets, realized by digitally controllable current sources, while termination is implemented with a digitally controllable resistor bank, providing tuning to compensate for process, voltage, temperature, and resistor process variations.

A modular and abuttable layout methodology is developed using common-centroid placement, symmetric routing, dummy devices, and continuous guard rings to reduce mismatch and layout-dependent effects. The complete transmitter occupies 176~$\mu$m $\times$ 82~$\mu$m. Pre- and post-extraction simulations are performed using a USB-IF reference channel with 21.6~dB insertion loss at the 5~GHz Nyquist frequency. Parasitic extraction significantly affects the output transition times, effective FFE tap weighting, and termination configuration, and tuning is required to maintain performance across different process corners.

The results show that the segmented architecture provides sufficient tuning range to recover USB~3.2 Gen~2-compatible operating points after physical implementation, but also demonstrate that the required equalization and termination codes can become dependent on process and operating conditions. %In particular, post-extraction equalization calibration cannot be represented by a single fixed mapping across all evaluated PVT corners. The work therefore demonstrates a complete schematic-to-layout methodology for a configurable 10~Gb/s transmitter while identifying automatic calibration, complete jitter/BER characterization, and silicon measurement as the principal steps required toward full experimental compliance validation.
In particular, the desired post-extraction equalization cannot be represented by a single fixed preset mapping across all evaluated PVT corners, which highlights the importance of calibration mechanisms and the versatility of the proposed architecture. The work therefore demonstrates a complete schematic-to-layout methodology for a configurable 10~Gb/s transmitter while identifying automatic calibration, complete jitter and BER characterization, and silicon measurement as the principal steps required toward full experimental compliance validation.

\vspace{10mm}
\textbf{Keywords:} High-speed transmitter, USB~3.2 Gen~2, Feed-Forward Equalization, Termination, Integrated-circuit design, Layout, Current-Mode Logic.

\acresetall %% to reset the acronym usage
